% Reusable layout snippets learned from stereochemistry v5.2
% These are design references, not a complete chapter source.

\usepackage[version=4]{mhchem}
\usepackage{chemfig}
\usepackage{tikz}
\usetikzlibrary{calc,arrows.meta,positioning}
\usepackage{booktabs,tabularx,array}

% Practice / answer dual-mode
\newif\ifanswers
\ifdefined\ANSWERS\answerstrue\else\answersfalse\fi

\newcommand{\A}[2][2.3cm]{%
  \ifanswers #2\else\rule{#1}{0.4pt}\fi
}

% Fischer projection: top, bottom, left, right
\newcommand{\Fischer}[4]{%
\begin{tikzpicture}[baseline=-0.55ex,scale=.58]
  \draw[line width=.72pt] (0,-1.15)--(0,1.15);
  \draw[line width=.72pt] (-1.15,0)--(1.15,0);
  \node[above] at (0,1.15) {#1};
  \node[below] at (0,-1.15) {#2};
  \node[left] at (-1.15,0) {#3};
  \node[right] at (1.15,0) {#4};
\end{tikzpicture}}

% Two-center Fischer
\newcommand{\FischerTwo}[6]{%
\begin{tikzpicture}[baseline=-0.5ex,scale=.54]
  \draw[line width=.72pt] (0,-1.85)--(0,1.85);
  \draw[line width=.72pt] (-1.0,.62)--(1.0,.62);
  \draw[line width=.72pt] (-1.0,-.62)--(1.0,-.62);
  \node[above] at (0,1.85) {#1};
  \node[left] at (-1.0,.62) {#2};
  \node[right] at (1.0,.62) {#3};
  \node[left] at (-1.0,-.62) {#4};
  \node[right] at (1.0,-.62) {#5};
  \node[below] at (0,-1.85) {#6};
\end{tikzpicture}}

% Newman: front top/left/right, back upper-right/upper-left/bottom
\newcommand{\Newman}[6]{%
\begin{tikzpicture}[baseline=-0.55ex,scale=.62]
  \draw[line width=.8pt] (0,0) circle (.47);
  \fill (0,0) circle (1.3pt);
  \foreach \ang/\lab in {90/#1,210/#2,330/#3}{
    \draw[line width=.72pt] (0,0)--(\ang:.94);
    \node at (\ang:1.20){\lab};
  }
  \foreach \ang/\lab in {30/#4,150/#5,270/#6}{
    \draw[line width=.72pt] (\ang:.47)--(\ang:.99);
    \node at (\ang:1.22){\lab};
  }
\end{tikzpicture}}

% Recommended pattern for structure -> CIP -> R/S -> full name
\begin{tabularx}{\linewidth}{>{\centering\arraybackslash}m{3.0cm} X c X}
\toprule
结构 & 优先顺序 & 构型 & 完整名称 / 结论\\
\midrule
\Fischer{\ce{CH3}}{\ce{C2H5}}{H}{OH}
& \A[4cm]{\ce{OH > C2H5 > CH3 > H}}
& \A[1cm]{R}
& \A[4cm]{(R)-丁-2-醇}\\
\bottomrule
\end{tabularx}

% Important layout lesson for four-option structure questions:
% If a 4-column layout makes labels collide, use a 2x2 grid or fixed-width minipages.
\begin{tabular}{cc}
\begin{minipage}{.45\linewidth}\centering A\\[2mm] <structure A>\end{minipage}
&
\begin{minipage}{.45\linewidth}\centering B\\[2mm] <structure B>\end{minipage}
\\[5mm]
\begin{minipage}{.45\linewidth}\centering C\\[2mm] <structure C>\end{minipage}
&
\begin{minipage}{.45\linewidth}\centering D\\[2mm] <structure D>\end{minipage}
\end{tabular}

% Rule: render the final PDF and inspect visually.
% A successful compile does NOT guarantee a readable chemical structure.
